% ============================================================
%  Exercices — Chapitre 00
%  Analyse dimensionnelle et remise à niveau
%  VERSION ENRICHIE : 50 exercices organisés en 10 thèmes (5-6 exos/thème)
% ============================================================

\section*{Chapitre 00 — Exercices (50 exercices)}

% ============================================================
\subsection*{Thème A — Vérification d'homogénéité (6 exercices)}
% ============================================================

\textbf{Niveau facile à moyen}

\subsubsection*{A1}  Perte de charge : $\Delta P = \lambda \dfrac{\ell}{d} \dfrac{\rho v^2}{2}$. Vérifier.

\subsubsection*{A2}  Bilan thermique : $\dot{Q} = h A \Delta T$. Vérifier.

\subsubsection*{A3}  Vitesse du son : $c = \sqrt{\dfrac{\gamma P}{\rho}}$. Vérifier.

\subsubsection*{A4}  Puissance hydraulique : $P_{\text{hyd}} = \dot{V} \times \Delta P$. Vérifier.

\subsubsection*{A5}  Nombre de Reynolds (sans dimension) : $\mathrm{Re} = \dfrac{\rho v d}{\mu}$. Vérifier.

\subsubsection*{A6}  Énergie cinétique volumique : $\rho v^2 = \Delta P$. Vérifier. Signification ?

% ============================================================
\subsection*{Thème B — Conversions simples (6 exercices)}
% ============================================================

\textbf{Niveau facile}

\subsubsection*{B1}  Convertir $P = 18~\mathrm{kW}$ en W, kWh/h, cal/s.

\subsubsection*{B2}  Convertir $\dot{V} = 720~\mathrm{m^3/h}$ en m³/s, L/min, L/s.

\subsubsection*{B3}  Convertir $\Delta P = 3{,}2~\mathrm{bar}$ en Pa, mH₂O, kPa.

\subsubsection*{B4}  Convertir $T = 75~\mathrm{°C}$ en K et °F. Différence T absolue vs Δ T ?

\subsubsection*{B5}  Convertir $\dot{m} = 2{,}5~\mathrm{kg/s}$ en g/s, tonnes/h.

\subsubsection*{B6}  Convertir $E = 125~\mathrm{kWh}$ en J, MJ, cal.

% ============================================================
\subsection*{Thème C — Conversions composées (6 exercices)}
% ============================================================

\textbf{Niveau facile à moyen}

\subsubsection*{C1}  Eau $\dot{V} = 2{,}4~\mathrm{m^3/h}$ (ρ = 1000 kg/m³). Convertir kg/s, tonnes/h.

\subsubsection*{C2}  Flux linéique tube $q_\ell = 45~\mathrm{W/m}$. Exprimer kW/(km), W/(100 m).

\subsubsection*{C3}  Flux 12 kW sur 8 m². Densité flux W/m², kW/m², W/cm² ?

\subsubsection*{C4}  Viscosité cinématique air : μ = 1,8×10⁻⁵ Pa·s, ρ = 1,2 kg/m³. Calculer ν en m²/s, cSt.

\subsubsection*{C5}  Air humide : $\dot{V} = 500~\mathrm{L/min}$ (ρ = 1,2), $h_v = 8~\mathrm{g/kg}$. Débit air m³/h, eau kg/h ?

\subsubsection*{C6}  Pompe : 15 kW absorbé, 12,3 kW hydraulique. Puissances W, kJ/h. Rendement % ?

% ============================================================
\subsection*{Thème D — Résolution d'équations linéaires (6 exercices)}
% ============================================================

\textbf{Niveau facile}

\subsubsection*{D1}  $\dot{Q} = \dot{m} c_p (T_{\text{in}} - T_{\text{out}})$. Exprimer $T_{\text{out}}$, $\dot{m}$, Δ T.

\subsubsection*{D2}  $\dot{Q} = h A (T_{\text{paroi}} - T_{\text{fluide}})$. Exprimer $T_{\text{paroi}}$, $h$, $A$.

\subsubsection*{D3}  Bifurcation : $\dot{V}_{\text{tot}} = 100~\mathrm{L/min}$, $\dot{V}_1 = 35~\mathrm{L/min}$. Trouver $\dot{V}_2$.

\subsubsection*{D4}  $\dot{Q} = k \Delta T$. Si $\dot{Q} = 8~\mathrm{kW}$, $\Delta T = 20~\mathrm{K}$, calculer $k$.

\subsubsection*{D5}  $\dot{m} = \rho \dot{V}$. Eau : ρ = 998 kg/m³, $\dot{V} = 0{,}05~\mathrm{m^3/s}$. Calculer $\dot{m}$. Exprimer $\dot{V}$.

\subsubsection*{D6}  $\Delta P = \lambda \dfrac{\ell}{d} \dfrac{\rho v^2}{2}$. Exprimer $v$, puis $\lambda$.

% ============================================================
\subsection*{Thème E — Équations non-linéaires (6 exercices)}
% ============================================================

\textbf{Niveau moyen}

\subsubsection*{E1}  $\dot{V} = v \pi \dfrac{d^2}{4}$. Exprimer $d$. Si $v = 2~\mathrm{m/s}$, $\dot{V} = 0{,}05~\mathrm{m^3/s}$, calculer.

\subsubsection*{E2}  $c = \sqrt{\gamma R T}$. Exprimer $T$ et $\gamma$. Vérifier air : γ = 1,4, R = 287, c ≈ 340 m/s.

\subsubsection*{E3}  $\mathrm{Re} = \dfrac{\rho v d}{\mu}$. Exprimer $v$, $d$. Re = 50000, ρ = 1000, μ = 10⁻³, d = 0,05 → $v$ ?

\subsubsection*{E4}  $P = \dot{V} \Delta P = v A \Delta P$. Exprimer $v$ en $P$, $A$, $\Delta P$.

\subsubsection*{E5}  Annulaire : $\dot{V} = v \pi (R_e^2 - R_i^2)$. Exprimer $R_e$.

\subsubsection*{E6}  $\varepsilon = 1 - \exp(-\mathrm{NUT})$. Exprimer $\mathrm{NUT}$. Calculer pour $\varepsilon = 0{,}7$.

% ============================================================
\subsection*{Thème F — Débits massique/volumique (6 exercices)}
% ============================================================

\textbf{Niveau facile à moyen}

\subsubsection*{F1}  Pompe $Q = 2{,}1~\mathrm{m^3/h}$. Convertir L/min, m³/s, L/s.

\subsubsection*{F2}  Extracteur $\dot{V} = 780~\mathrm{m^3/h}$ air (ρ = 1,2). Convertir m³/s, L/min. Débit massique kg/s, t/h.

\subsubsection*{F3}  Compresseur 45 m³/h (1 atm). L/min ? ρ = 1,2 → $\dot{m}$ kg/s ?

\subsubsection*{F4}  Frigorifique $\dot{m} = 0{,}15~\mathrm{kg/s}$ (ρ = 1300 kg/m³). $\dot{V}$ m³/h, L/min, cm³/s ?

\subsubsection*{F5}  Air humide $\dot{V} = 120~\mathrm{m^3/h}$ (ρ = 1,18), $h_v = 10~\mathrm{g/kg}$. Débit sec, eau kg/h ?

\subsubsection*{F6}  Air réchauffe 20°C (ρ₁ = 1,204), 500 m³/h → 60°C (ρ₂ = 1,067). Nouveau $\dot{V}_2$ ? $\dot{m}$ change ?

% ============================================================
\subsection*{Thème G — Analyse dimensionnelle par exposants (6 exercices)}
% ============================================================

\textbf{Niveau moyen à difficile}

\subsubsection*{G1}  $\Delta P = k \rho^\alpha v^\beta \ell^\gamma d^\delta$. Trouver α, β, γ, δ.

\subsubsection*{G2}  $F = k \rho^\alpha v^\beta A^\gamma$ (traînée). Déterminer exposants.

\subsubsection*{G3}  $\dot{V} = k A^\alpha \sqrt{g^\beta h^\gamma}$ (orifice). Trouver α, β, γ.

\subsubsection*{G4}  $P = k \rho^\alpha g^\beta h^\gamma \dot{V}^\delta$ (pompe). Déterminer.

\subsubsection*{G5}  Vérifier $[\nu] = [\mathrm{m^2/s}]$ pour $\nu = \mu/\rho$ (Pa·s / kg/m³).

\subsubsection*{G6}  Strouhal adimensionnel : $\mathrm{St} = \dfrac{f d}{v}$. Vérifier (f Hz, d m, v m/s).

% ============================================================
\subsection*{Thème H — Ordre de grandeur (5 exercices)}
% ============================================================

\textbf{Niveau facile}

\subsubsection*{H1}  Classer puissances : 10 W (LED) ; 1500 W (radiateur) ; 3 kW (clim) ; 100 kW (chaudière). 
Puissances lycée ?

\subsubsection*{H2}  Vitesse air 15 m/s dans gaine. Plausible ? (Norme confort 3–6 m/s, max 8–12 m/s).

\subsubsection*{H3}  ECS 5 m³/h pour 20 pers. Volume/pers/h ? Plausible douche 10 min ?

\subsubsection*{H4}  Échangeur Δ T = −3 K. Possible ?  Pourquoi ?

\subsubsection*{H5}  Paroi intérieure −15°C (extérieur 5°C). Plausible ? Conséquences ?

% ============================================================
\subsection*{Thème I — Interpolation/abaques (5 exercices)}
% ============================================================

\textbf{Niveau facile à moyen}

\subsubsection*{I1}  Tableau eau (ρ, $c_p$, μ) 0, 20, 40, 60, 80, 100°C. Interpoler 30°C, 50°C, 55°C.

\subsubsection*{I2}  Tableau air (ρ, ν) 0, 10, 20, 30°C. Interpoler 25°C.

\subsubsection*{I3}  Moody (Re = 1000, 2000, 5000, 10000) λ. Interpoler log-log Re = 3000.

\subsubsection*{I4}  Viscosité fluide (25, 35, 45°C). Estimé 40°C ? Croît ou décroît T ?

\subsubsection*{I5}  Enthalpie vs T (10°C : 130 kJ/kg ; 30°C : 145). Interpoler 20°C. Linéaire ?

% ============================================================
\subsection*{Thème J — Synthèse complète (5 exercices)}
% ============================================================

\textbf{Niveau moyen à difficile}

\subsubsection*{J1 — Pompe}
$Q = 1{,}2~\mathrm{m^3/h}$, $\Delta P = 25\,000~\mathrm{Pa}$. 
a) Convertir. b) Puissance hydraulique. c) η = 70 % → puissance moteur. d) Plausible ?

\subsubsection*{J2 — Radiateur}
12 kW, eau 70→55°C ($c_p = 4{,}18~\mathrm{kJ/(kg \cdot K)}$). 
a) Débit massique. b) kg/h, L/min (ρ = 985). c) Vérifier bilan. d) Si Δ T = 5 K ?

\subsubsection*{J3 — Conduit}
$\dot{V} = 400~\mathrm{m^3/h}$, d = 0,2 m, ℓ = 50 m, λ = 0,025. 
a) Débit convert. b) Vitesse. c) Reynolds (ρ = 1,2, μ = 1,8×10⁻⁵). d) Darcy. e) Pa, bar, mH₂O.

\subsubsection*{J4 — Récupération}
Effluent 30→20°C, $\dot{m} = 5~\mathrm{kg/s}$ ($c_p = 4{,}18~\mathrm{kJ/(kg \cdot K)}$). 
a) Puissance kW. b) Débit m³/h. c) Énergie 24h kWh. d) Coût 0,12 €/kWh/jour.

\subsubsection*{J5 — Frigorifique R410A}
$\dot{m} = 0{,}08~\mathrm{kg/s}$, ρ = 850 kg/m³, Δh = 180 kJ/kg. 
a) $\dot{V}$ m³/s, L/min, m³/h. b) Puissance frigorifique kW. c) Moteur 15 kW → COP. 
d) COP = 3 plausible ? e) Vérifier dimensions.

\newpage

% ============================================================
%  CORRIGÉS — SYNTHÈSE
% ============================================================

\section*{Corrigés — Chapitre 00}

\subsection*{Thème A — Corrigés}

\subsubsection*{A1–A6}  ✓ Toutes homogènes. A6 = énergie cinétique volumique.

\subsection*{Thème B — Corrigés (résumé)}

B1: 18 000 W = 18 kWh/h ≈ 4 300 cal/s | 
B2: 0,2 m³/s = 12 000 L/min |
B3: 320 000 Pa = 32 mH₂O = 320 kPa |
B4: 348 K = 167°F |
B5: 2 500 g/s = 9 t/h |
B6: 4,5×10⁸ J = 450 MJ

\subsection*{Thème C — Corrigés (résumé)}

C1: 2,4×10⁶ kg/h = 2 400 t/h |
C2: 4,5 kW/(100 m) |
C3: 1 500 W/m² = 1,5 kW/m² |
C4: ν = 1,5×10⁻⁵ m²/s = 0,015 cSt |
C5: 30 m³/h air | 0,3 kg/h eau |
C6: η = 82 % ✓

\subsection*{Thèmes D, E, F, G, H, I, J — Corrigés}

\textit{Voir document séparé « Corrigés-Complets-Chapitre-00.pdf »}

\textit{Clés pédagogiques :}
\begin{itemize}
  \item \textbf{D} : Isoler variable, vérifier dimensions
  \item \textbf{E} : Racines/puissances → 2 solutions possibles
  \item \textbf{F} : Relation ρ, $\dot{m}$, $\dot{V}$ fondamentale
  \item \textbf{G} : Système 3 équations / 4 inconnues
  \item \textbf{H} : Valider ordre de grandeur toujours
  \item \textbf{I} : Interpolation linéaire, log si nécessaire
  \item \textbf{J} : Synthèse 3–4 techniques
\end{itemize}

\end{document}
